% !PRINTED 288
\weberpage{288}
%<<CONT
and we shall further assume that it passes through none of the
points of intersection of $\varphi$ and $\psi$.

We give these three functions a definite cyclic order
\[
f,\ \varphi,\ \psi,
\]
so that $f$ is again to follow $\psi$, and we associate with this
order a definite direction of advance on each of the three curves,
say the positive one. We shall stipulate that with the reversed
order
\[
f,\ \psi,\ \varphi
\]
the opposite, that is the negative, direction of advance shall be
associated on each of the three curves.

We now traverse, in the first order, the curve $f$ in the
positive sense, and attend to the points of intersection of $f$
and $\varphi$. Such a point shall be called an exit point $A(f; \varphi, \psi)$

\weberfigureimage{weber-b1-fig08.png}{0.48}

if the positive direction of $f$ leads from the interior of the
inner region $(\varphi, \psi)$ into the exterior, and an entry point
$E(f; \varphi, \psi)$ if it leads from the exterior into the interior.

In Fig.~8 the points $a$ are the exit points, the point
$e$ an entry point.

The total number of entry and exit points agrees with the
number of points of intersection of $f$ and $\varphi$, and is therefore, since
both curves are closed, an \gesperrt{even number}. If $a$ is the number
of points $A(f; \varphi, \psi)$, $e$ the number of points $E(f; \varphi, \psi)$, then
$e - a$ is likewise an even number, and
\[
\tfrac{1}{2}(e - a) = k
\]
is an integer, which after \gesperrt{Kronecker} is called the \gesperrt{characteristic
of the system of functions} $f$, $\varphi$, $\psi$. (In the case of Fig.~8
it is equal to $-1$).

If we permute the functions $f$, $\varphi$, $\psi$ cyclically, we
obtain three determinations of the characteristic, and if
we reverse the order and thereby, in accordance with the convention
adopted, also replace the positive directions by the negative
ones, three more. These determinations are, if we now use
the symbols $A(f; \varphi, \psi)$, $E(f; \varphi, \psi)$ at the same time as a designation
for the numbers of the points concerned,
%CONT>>
